% Recuperación documental para el artículo V; RESULTADO_RECUPERADO.
% Propietario: /Users/ruben/Documents/New project/output/REVISION_EDITORIAL_HMT_MD_20260905/01_FUENTES/INTEGRAL/tree/New project/output/ENSAMBLADO_CAUSAL_RESTAURADO_HMT_MD_20260905/fuente/manuscrito/sections/md/delta_127/04b_completaciones_cuanticas.tex
% Líneas de origen: 1--334.
% REV02: construcción algebraica; la aplicación estadística se reúne en 40.
% Correspondencia editorial: gestion/CONSOLIDACION_PAULI_FOCK_REV02.json.
\section{Construcción algebraica de las completaciones cuánticas sobre APP}
\label{v:rec:chap:completaciones-cuanticas}
\label{v:sec:completaciones-algebraicas}

APP y TPK proporcionan estados orientados, transporte y paridad de hoja.
Su realización lineal permite construir las álgebras simétrica y exterior,
con sus operadores de creación, aniquilación y ocupación. Esta sección
establece esas construcciones algebraicas, el sector celular de campos y
el acoplamiento reversible de medición. La aplicación de los espectros de ocupación al equilibrio se desarrolla en la sección~\ref{v:rec:ocupacion:section:02}.
Cada tramo conserva las condiciones de la representación que utiliza.

\subsection{Espacio lineal de estados asociado a una realización de partícula}
\label{v:rec:completaciones:section:02}

Sea \(S\) un conjunto finito de clases de rutas TPK a la profundidad elegida
y sea
\[
 \mathcal H_1=\ell^2(S)
\]
el espacio de Hilbert con base ortonormal \((e_s)_{s\in S}\). Una elección de
pesos positivos \(E_s\) define el operador diagonal
\[
 He_s=E_se_s.
\]
La elección de \(S\), de su cociente por simetrías y de los pesos forma parte
del modelo. Una vez fijada, las completaciones multipartícula son canónicas.

\subsection{Completaciones simétrica y exterior}
\label{v:rec:completaciones:section:03}

Definimos
\[
 \mathcal F_+(\mathcal H_1)=
 \bigoplus_{n\geq0}\operatorname{Sym}^n\mathcal H_1,
 \qquad
 \mathcal F_-(\mathcal H_1)=
 \bigoplus_{n\geq0}\bigwedge^n\mathcal H_1.
\]
La primera es la completación bosónica; la segunda, la fermiónica. En ambas,
el sumando de grado cero es el espacio del vacío. Las potencias exteriores
se dotan del producto escalar para el cual las cuñas ordenadas de vectores
de una base ortonormal son ortonormales; las potencias simétricas se utilizan
con la normalización usual de los operadores de creación y aniquilación.

\subsection{Graduación de hoja y álgebra de intercambio}
\label{v:rec:completaciones:section:04}

Supongamos que la hoja orientada define un carácter aditivo
\[
 p:\operatorname{Mor}(\mathsf P_{\rm TPK})\longrightarrow\mathbb Z/2\mathbb Z,
 \qquad p(\gamma\eta)=p(\gamma)+p(\eta).
\]
Sobre el espacio lineal graduado correspondiente se introduce la simetría
\[
 \tau(v\otimes w)=(-1)^{p(v)p(w)}w\otimes v.
\]

\begin{theorem}[Completación estadística determinada por la paridad]
\label{v:rec:completaciones:statement:02}
El cociente del álgebra tensorial por las relaciones
\[
 v\otimes w-(-1)^{p(v)p(w)}w\otimes v
\]
es simétrico sobre el sector par y exterior sobre el sector impar. Dos estados
impares idénticos tienen producto nulo; los estados pares admiten ocupación
arbitraria.
\end{theorem}

\begin{proof}
Si \(p(v)=p(w)=0\), la relación es \(vw=wv\). Si
\(p(v)=p(w)=1\), es \(vw=-wv\); tomando \(v=w\) y trabajando en
característica cero se obtiene \(v^2=0\). Los sectores mixtos conmutan sin
signo adicional. Esta es la propiedad universal del álgebra simétrica
graduada.
\end{proof}

\subsection{Operadores canónicos e idempotencia de la ocupación}
\label{v:pauli:car-construccion}

El carácter de paridad determina la completación; el producto de esa
completación determina a continuación el álgebra de operadores. En el
sector exterior, la creación \(a_s^\dagger\) es el producto exterior por
\(e_s\), y \(a_s\) es su adjunto.

\begin{theorem}[Principio de exclusión de Pauli en la completación exterior]
\label{v:rec:completaciones:statement:01}
\label{v:rec:thm:principio-exclusion-pauli}
Para todo modo \(s\in S\), los operadores de creación y aniquilación de la
completación exterior satisfacen
\begin{equation}
 \{a_s,a_t^\dagger\}=\delta_{st}I,
 \qquad
 \{a_s,a_t\}=\{a_s^\dagger,a_t^\dagger\}=0.
 \label{v:pauli:eq:car}
\end{equation}
En particular, el operador de ocupación
\(N_s=a_s^\dagger a_s\) es un proyector:
\begin{equation}
 N_s^2=N_s,
 \qquad \Spec(N_s)\subseteq\{0,1\}.
 \label{v:pauli:eq:proyeccion-ocupacion}
\end{equation}
Los autovalores \(0\) y \(1\) corresponden, respectivamente, a la ausencia
y a la presencia del modo \(s\).
\end{theorem}

\begin{proof}
Sobre una cuña ordenada, \(a_s\) suprime el factor \(e_s\), cuando está
presente, con el signo de su posición; si está ausente, da cero. Componer
esta operación con el producto exterior demuestra las relaciones
\eqref{v:pauli:eq:car}: para \(s=t\), las dos composiciones cubren los casos
complementarios de ausencia y presencia; para \(s\ne t\), los signos de
intercambio se cancelan. La antisimetría implica
\(e_s\wedge e_s=0\) y
\(a_s^2=(a_s^\dagger)^2=0\). Entonces
\begin{equation}
 N_s^2=a_s^\dagger a_s a_s^\dagger a_s
 =a_s^\dagger(I-a_s^\dagger a_s)a_s
 =a_s^\dagger a_s=N_s.
 \label{v:pauli:eq:prueba-idempotencia}
\end{equation}
El operador es autoadjunto, y todo autovalor de un idempotente satisface
\(\lambda^2=\lambda\). Por tanto, su espectro está contenido en
\(\{0,1\}\). El vacío y el estado \(e_s\) realizan ambos valores
\cite{Pauli1925,Dirac1926}.
\end{proof}

En el sector simétrico, los operadores bosónicos \(b_s^\dagger,b_s\)
actúan sobre la base normalizada de ocupaciones por los factores
\(\sqrt{m_s+1}\) y \(\sqrt{m_s}\), respectivamente. De sus composiciones
se obtienen
\begin{equation}
 [b_s,b_t^\dagger]=\delta_{st}I,
 \qquad [b_s,b_t]=[b_s^\dagger,b_t^\dagger]=0.
 \label{v:pauli:eq:ccr}
\end{equation}
Estas identidades se entienden sobre el dominio algebraico de partículas
finitas, estable bajo los operadores considerados.

El teorema~\ref{v:rec:completaciones:statement:02} convierte una paridad composicional TPK en una regla exacta de
ocupación. Para obtener el teorema físico de espín--estadística hay que
construir además una teoría local lorentziana positiva y demostrar que la
paridad de hoja coincide con el espín bajo rotaciones. La igualdad no se
deduce únicamente de que ambos objetos tomen valores módulo dos.
